\section{Bài báo giải được gì, và để lại gì}
\label{sec:ch07-con-lai-gi}

\index{Vaswani 2017!bài giải được gì}%
Phát biểu cho gọn: bài bỏ được ràng buộc rằng một câu \(n\) token phải là \(n\)
lượt tính nối đuôi nhau, và bỏ bằng cách thay recurrence bằng attention của câu
với chính nó. Cột \code{Sequential} của bảng 1 đi từ \(O(n)\) xuống \(O(1)\), và
đó là toàn bộ chuyện.

Ba món nợ mà chương~\ref{ch:encoder-decoder} để lại thì nay trả xong cả ba.
Chương~\ref{ch:mot-vector-cho-moi-tu} trả hai món đầu, encoder không còn vứt bỏ
trạng thái và decoder hỏi lại được câu nguồn. Món thứ ba, tính tuần tự, là món
chương này trả.

\subsection*{Con số của bài, và một chỗ phải nói rõ}

\index{Vaswani 2017!41.0 hay 41.8}%
Bảng 2 báo BLEU 27.3 cho model base và 28.4 cho model big trên WMT 2014
Anh-Đức, và 38.1 cho base trên Anh-Pháp. Model big huấn luyện 3.5 ngày trên tám
GPU P100.

Còn con số Anh-Pháp của model big thì tùy bạn cầm bản nào.

\begin{itemize}
  \item Bản kỷ yếu NIPS 2017 ghi \textbf{41.0} ở tóm tắt, ở bảng 2, và ở thân
    bài mục 6.1. Ba chỗ khớp nhau.
  \item Bản arXiv v7, đề tháng 8 năm 2023, ghi \textbf{41.8} ở tóm tắt và ở
    bảng 2, nhưng thân bài mục 6.1 của chính nó vẫn ghi \enquote{our big model
    achieves a BLEU score of 41.0}.
\end{itemize}

Tôi mở cả hai file ở trang 8 để chắc, vì đây đúng là loại chuyện mà đọc lướt sẽ
bỏ qua. Con số 41.8 mà phần lớn chỗ trích lại dùng chỉ có trong bảng và tóm tắt
của một bản sửa sáu năm sau hội nghị, và bản ấy tự mâu thuẫn với chính nó cách
đó hai đoạn.

Tôi không biết vì sao, và tôi không đoán. Cái rút ra được thì không cần biết vì
sao: một bài báo không phải một vật cố định, và \enquote{Vaswani và cộng sự báo
41.8} là một câu chưa đủ. Câu đủ phải nói bản nào. Cùng bài học mà
mục~\ref{sec:ch06-ai-do} rút ra từ chuyện đếm trang PDF, một tầng cao hơn: thứ
bạn đang đọc cũng là một phép đo, và nó cũng sai được.

\subsection*{Cái bài tự nêu là còn thiếu}

\index{Vaswani 2017!phần kết luận nêu gì}%
Phần kết luận nêu ba hướng, và câu cuối là câu đáng đọc nhất:

\begin{quote}
\enquote{We plan to extend the Transformer to problems involving input and
output modalities other than text and to investigate local, restricted
attention mechanisms to efficiently handle large inputs and outputs such as
images, audio and video. Making generation less sequential is another research
goals of ours.}
\end{quote}

Câu thứ nhất trỏ thẳng vào chương~\ref{ch:vit}: đưa kiến trúc này sang ảnh. Vế
sau của cùng câu ấy, \enquote{local, restricted attention}, là dòng nghiên cứu
mà chương~\ref{ch:chi-phi-bac-hai} kể lại, kể cả chỗ nó không sống sót.

Câu thứ hai là câu tôi muốn bạn dừng lại, và nó cũng là câu sai ngữ pháp trong
chính bài, ở cả ba bản tôi xem. \enquote{Making generation less sequential}.
Bài vừa bỏ tính tuần tự đi xong, và ở dòng cuối bài tự nói rằng nó vẫn còn.

Vì nó còn thật, và đây là chỗ dễ nhầm nhất về kiến trúc này. Cái được song song
hóa là \emph{huấn luyện}, nơi cả câu đích có sẵn nên mọi vị trí chạy một lượt.
Lúc sinh câu thì token \(i+1\) chưa tồn tại trước khi token \(i\) sinh ra, nên
decoder vẫn đi từng bước. Tệ hơn: không còn trạng thái để mang theo, nên mỗi
bước phải chạy lại decoder trên toàn bộ prefix. Hàm giải mã trong repo làm đúng
thế, và tôi để nguyên chứ không tối ưu, vì cái làm nó hết lãng phí là bộ nhớ
đệm khóa và giá trị, tức chuyện của chương~\ref{ch:thuc-te-trien-khai}.

\subsection*{Cái bài không nêu, và chương này đo được}

\index{Vaswani 2017!hai chỗ chương này đo được}%
Hai chuyện, cả hai từ số đo của chương chứ không từ bài.

\textbf{Thứ nhất, thiên kiến quy nạp.} Mục~\ref{sec:ch07-corpus} đo được rằng ở
đúng recipe chương trước, kiến trúc này được 0.4700 so với 1.0000, và cần ba
lần số epoch mới tới cùng chỗ dù mang ít tham số hơn. Kiểu lỗi còn lại thì đều
một cách đáng chú ý: đúng từ, sai cụm. Quan hệ \enquote{đứng ngay trước} là thứ
mạng hồi quy có sẵn và self-attention phải học. Bảng 1 tính chi phí mỗi tầng và
độ dài đường đi; nó không có cột nào cho chuyện phải học bao nhiêu mới biết hai
vị trí cạnh nhau. Chương~\ref{ch:bias-quy-nap} là chương ấy, và
chương~\ref{ch:vit} là chỗ cùng câu hỏi quay lại ở quy mô lớn, nơi câu trả lời
là dữ liệu đủ nhiều thì mua được thiên kiến.

\textbf{Thứ hai, cái giá bậc hai giờ nằm ở mọi chỗ.} Bahdanau tự gọi tích
\(T_x \times T_y\) là một \enquote{drawback} rồi gạt đi vì câu dịch chỉ dài 15
tới 40 từ, và mục~\ref{sec:ch06-con-lai-gi} đã đọc kỹ ba câu ấy. Lý do gạt bỏ
ấy hết đúng ở đây, theo hai hướng cùng lúc. Cơ chế không còn chạy một lần giữa
hai câu, nó chạy trong mỗi tầng và ở cả ba chỗ của
mục~\ref{sec:ch07-tu-nhin}, kể cả nối câu nguồn với chính nó, nên tích thành
\(n^2\) nhân với số tầng. Và \(n\) thôi không còn là 15 tới 40.

Ô \(O(n^2 \cdot d)\) ở góc bảng 1 là ô bài in ra mà không coi là vấn đề, vì ở
\(n\) nhỏ hơn \(d\) thì nó nhỏ hơn \(O(n \cdot d^2)\) của mạng hồi quy, và bài
nói đúng chuyện đó. Chương~\ref{ch:chi-phi-va-song-song} đo lại ô ấy bằng FLOP,
bằng bộ nhớ và bằng đồng hồ; chương~\ref{ch:chi-phi-bac-hai} là chương của mười
năm sau, khi \(n\) đã vượt \(d\) từ lâu.

\subsection*{Chỗ chương sau bắt đầu}

Kiến trúc thì xong ở đây. Cái chưa xong là ba lựa chọn bài đưa ra mà không giải
thích: \eqref{eq:ch07-post-ln} viết layer norm sau phép cộng chứ không phải
trước, mục 5.3 có một lịch warmup 4000 bước, và mục 5.4 có label smoothing mà
bài thừa nhận là \enquote{hurts perplexity}. Đọc lướt thì cả ba là mẹo chỉnh
tay.

Chương này không có tư cách nói chúng là mẹo hay không, và chỗ ấy đáng nói
thẳng. Mọi mô hình trong mục~\ref{sec:ch07-corpus} chạy \emph{không} warmup,
\emph{không} label smoothing và dropout bằng 0, và chúng vẫn đạt 0.9967. Nên ở
quy mô này, ba thứ ấy không cần thiết, và một chương kết luận ngược lại từ đây
là một chương đọc quá tay bảng của chính nó.

Cái làm chúng đáng đo nằm ở chỗ khác: bài chạy \(\dmodel = 512\), sáu tầng mỗi
phía, 300000 bước, và ở quy mô ấy thứ tự của layer norm đổi hẳn hình dạng
gradient lúc khởi động. Chương~\ref{ch:chi-phi-va-song-song} là chỗ đo cả ba ở
quy mô mà chúng bắt đầu có nghĩa, và câu hỏi nó phải trả lời là một câu hỏi về
quy mô chứ không phải một câu hỏi về ba mẹo.
